\section{Conclusions}
The following conclusions are drawn based on our investigation of partitioning and density of equilibrium phases in \ce{CO2}-water and \ce{CO2}-brine mixtures from CPA/eCPA Eos and GEMC simulations:

\begin{enumerate}
    \item We offer a new set of parameters of the CPA EoS to describe more accurately the phase diagram of \ce{CO2}-water at high temperatures using two temperature-dependent interaction parameters. The model accurately describes partitioning in the water-rich phase in the entire range of temperature (283 K$ \leq T \leq$ 623 K) and in the \ce{CO2}-rich phase for $T$ > 304 K;
    \item \textcolor{red}{We parametrized the eCPA EoS to model the \ce{CO2} solubility in NaCl brine at low to high temperatures and from dilute to concentrated salt mixtures. Compared to the past model, the improvement is 74 \% in the average deviation of \ce{CO2} solubility in brine, which is obtained by accounting for the \ce{CO2}-\ce{H2O} cross-association and including temperature-dependency on the parameters for \ce{CO2}-salt and water-\ce{CO2} interactions. The ternary mixture is modeled with 4 binary parameters, which depend on temperature by a quadratic polynomial};
    \item We investigate the \ce{CO2}-water and \ce{CO2}-brine mixture density at equilibrium and below the saturation. Good agreement with experiments is obtained with deviations of about 1\% for the aqueous phase. \ce{CO2} dissolution may increase or decrease the brine density depending on salinity and temperature;
    \item From GEMC simulations, we show that SPCE-water and EPM2-\ce{CO2} are a good combination for the high-temperature phase equilibria when unlike parameters are chosen properly. At such conditions, thermal effects overcome the SPCE overpolarization, and the water solubility in the \ce{CO2}-rich phase is computed without further corrections;
    \item We compare the association from EoS and molecular simulations. Similar water self-association and \ce{CO2}-water cross-association are obtained by the two different techniques. The significance of \ce{CO2} cross-association is established from a molecular structure perspective.
\end{enumerate}

With the modeling tools and analysis developed in our work, we hope to improve the thermodynamics simulation of subsurface formations. In particular, we highlight conditions in geothermal formations, which have high temperatures.
